% Example operating points, computed from the analytic model equations.
% -------------------------------------------------------------------------
% MTPC point.
\pgfmathsetmacro{\cA}{3.0}
\pgfmathsetmacro{\yA}{0.52}
\pgfmathsetmacro{\xA}{xmtpcnorm(\yA)}
\pgfmathsetmacro{\rA}{sqrt(\xA*\xA+\yA*\yA)}
\pgfmathsetmacro{\TA}{torque(\Imax*\xA,\Imax*\yA)}

% Feasible FW point: on voltage ellipse, inside current circle, left of MTPC.
\pgfmathsetmacro{\cB}{1.8}
\pgfmathsetmacro{\xB}{-0.65}
\pgfmathsetmacro{\yB}{yvoltnorm(\xB,\cB)}
\pgfmathsetmacro{\TB}{torque(\Imax*\xB,\Imax*\yB)}

% MC point: analytic intersection of current circle and voltage ellipse.
\pgfmathsetmacro{\cC}{2.2}
\pgfmathsetmacro{\Aell}{\Ld*\Imax}
\pgfmathsetmacro{\Bell}{\Lq*\Imax}
\pgfmathsetmacro{\qaC}{\Aell*\Aell-\Bell*\Bell}
\pgfmathsetmacro{\qbC}{2*\psim*\Aell}
\pgfmathsetmacro{\qcC}{\psim*\psim+\Bell*\Bell-\cC*\cC}
\pgfmathsetmacro{\discC}{\qbC*\qbC-4*\qaC*\qcC}
\pgfmathsetmacro{\rootCone}{(-\qbC+sqrt(\discC))/(2*\qaC)}
\pgfmathsetmacro{\rootCtwo}{(-\qbC-sqrt(\discC))/(2*\qaC)}
\pgfmathsetmacro{\xC}{(\rootCone<0 && \rootCone>-1) ? \rootCone : \rootCtwo}
\pgfmathsetmacro{\yC}{ycurrentnorm(\xC)}
\pgfmathsetmacro{\TC}{torque(\Imax*\xC,\Imax*\yC)}

% Determine intersections for the feasible range and MTPV point.
\pgfmathsetmacro{\bestErrM}{1000}
\pgfmathsetmacro{\yM}{0.66}
\foreach \k in {0,...,1000}{%
        \pgfmathsetmacro{\ytmp}{\k/1000}%
        \pgfmathsetmacro{\errtmp}{abs(xcurrentnorm(\ytmp)-xmtpvnorm(\ytmp))}%
        \pgfmathparse{\errtmp < \bestErrM ? 1 : 0}%
        \ifnum\pgfmathresult=1
            \xdef\bestErrM{\errtmp}%
            \xdef\yM{\ytmp}%
        \fi
    }
\pgfmathsetmacro{\bestErrR}{1000}
\pgfmathsetmacro{\yR}{0.93}
\foreach \k in {0,...,1000}{%
        \pgfmathsetmacro{\ytmp}{\k/1000}%
        \pgfmathsetmacro{\errtmp}{abs(xcurrentnorm(\ytmp)-xmtpcnorm(\ytmp))}%
        \pgfmathparse{\errtmp < \bestErrR ? 1 : 0}%
        \ifnum\pgfmathresult=1
            \xdef\bestErrR{\errtmp}%
            \xdef\yR{\ytmp}%
        \fi
    }

% MTPV operating point: chosen on the current-circle/MTPV intersection.
\pgfmathsetmacro{\yD}{0.25}
\pgfmathsetmacro{\xD}{xmtpvnorm(\yD)}
\pgfmathsetmacro{\cD}{sqrt((\psim+\Ld*\Imax*\xD)^2+(\Lq*\Imax*\yD)^2)}
\pgfmathsetmacro{\TD}{torque(\Imax*\xD,\Imax*\yD)}

% plot configuration for the operating range slides
\pgfplotsset{
    opsaxis/.style={
            width=10.5cm,
            height=7.5cm,
            xmin=-1.05, xmax=0.05,
            ymin=0.0, ymax=1.10,
            axis lines=left,
            axis equal image,
            xlabel={$i_{\mathrm{s,d}}/i_{\mathrm{s,max}}$},
            ylabel={$i_{\mathrm{s,q}}/i_{\mathrm{s,max}}$},
            grid=major,
            tick label style={font=\small},
            label style={font=\small}
        },
    opsaxisgen/.style={
            width=10.5cm,
            height=8.0cm,
            xmin=-1.05, xmax=0.05,
            ymin=-1.05, ymax=1.05,
            axis lines=left,
            axis equal image,
            xlabel={$i_{\mathrm{s,d}}/i_{\mathrm{s,max}}$},
            ylabel={$i_{\mathrm{s,q}}/i_{\mathrm{s,max}}$},
            grid=major,
            tick label style={font=\small},
            label style={font=\small}
        }
}


% For speed-torque map slide

\pgfmathdeclarefunction{opsTSfluxnorm}{2}{%
    \pgfmathparse{sqrt((\psim+\Ld*\Imax*(#1))^2 + (\Lq*\Imax*(#2))^2)}%
}

% Nominal point: intersection of MTPC and current limit.
\pgfmathsetmacro{\opsTSxNom}{xmtpcnorm(\yR)}
\pgfmathsetmacro{\opsTSTNom}{torque(\Imax*\opsTSxNom,\Imax*\yR)}
\pgfmathsetmacro{\opsTScNom}{opsTSfluxnorm(\opsTSxNom,\yR)}

% MTPF/MTPV cut-in point: intersection of high-speed locus and current limit.
\pgfmathsetmacro{\opsTSxMTPFcut}{xcurrentnorm(\yM)}
\pgfmathsetmacro{\opsTSomegaMTPFcut}{\opsTScNom/opsTSfluxnorm(\opsTSxMTPFcut,\yM)}
\pgfmathsetmacro{\opsTSTorqueMTPFcut}{torque(\Imax*\opsTSxMTPFcut,\Imax*\yM)/\opsTSTNom}

% Alias names, in case existing slide text uses MTPV nomenclature.
\pgfmathsetmacro{\opsTSomegaMTPVcut}{\opsTSomegaMTPFcut}
\pgfmathsetmacro{\opsTSTorqueMTPVcut}{\opsTSTorqueMTPFcut}

% MTPC feasibility speed at zero torque, used to close the shaded regions.
\pgfmathsetmacro{\opsTSomegaMTPCzero}{\opsTScNom/opsTSfluxnorm(xmtpcnorm(0),0)}

% Select which already defined current-plane example shall be shown in the
% torque-speed map. Here: point B from the FW-region slide.
% To show A, C or D instead, replace (\cB,\TB) by (\cA,\TA), (\cC,\TC) or
% (\cD,\TD), respectively.
\pgfmathsetmacro{\opsTSomegaOp}{\opsTScNom/\cB}
\pgfmathsetmacro{\opsTSTorqueOp}{\TB/\opsTSTNom}

% Visible high-speed range and the corresponding endpoint of the MTPF curve.
\pgfmathsetmacro{\opsTSomegaMaxPlot}{3.08}
\pgfmathsetmacro{\opsTSbestErrEnd}{1000}
\pgfmathsetmacro{\opsTSyEnd}{0.20}
\foreach \k in {400,...,1000}{%
        \pgfmathsetmacro{\opsTSytmp}{\k*\yM/1000}%
        \pgfmathsetmacro{\opsTSomegaTmp}{\opsTScNom/opsTSfluxnorm(xmtpvnorm(\opsTSytmp),\opsTSytmp)}%
        \pgfmathsetmacro{\opsTSerrTmp}{abs(\opsTSomegaTmp-\opsTSomegaMaxPlot)}%
        \pgfmathparse{\opsTSerrTmp < \opsTSbestErrEnd ? 1 : 0}%
        \ifnum\pgfmathresult=1
            \xdef\opsTSbestErrEnd{\opsTSerrTmp}%
            \xdef\opsTSyEnd{\opsTSytmp}%
        \fi
    }

% Precomputed label anchor positions avoid arithmetic inside axis coordinates.
\pgfmathsetmacro{\opsTSomegaCutLabel}{\opsTSomegaMTPFcut+0.06}
\pgfmathsetmacro{\opsTSTorqueCutLabel}{\opsTSTorqueMTPFcut+0.04}
\pgfmathsetmacro{\opsTSomegaOpLabel}{\opsTSomegaOp+0.04}
\pgfmathsetmacro{\opsTSTorqueOpLabel}{\opsTSTorqueOp+0.02}

% Functions for plotting the contours parametrically by the normalized
% q-axis current y.
\pgfmathdeclarefunction{opsTSomeganormxy}{2}{%
    \pgfmathparse{\opsTScNom/opsTSfluxnorm(#1,#2)}%
}
\pgfmathdeclarefunction{opsTStorquenormxy}{2}{%
    \pgfmathparse{torque(\Imax*(#1),\Imax*(#2))/\opsTSTNom}%
}

\pgfplotsset{
    opsspeedaxis/.style={
            width=13.4cm,
            height=7.15cm,
            xmin=0, xmax=3.14,
            ymin=0, ymax=1.10,
            axis lines=left,
            axis line style={-Latex},
            xlabel={$\omega/\omega_{\mathrm{nom}}$},
            ylabel={$T^*/T_{\mathrm{max}}$},
            grid=major,
            major grid style={gray!35},
            tick label style={font=\small},
            label style={font=\small},
            clip=false,
            legend style={draw=none,fill=none,font=\small},
        }
}

%% For LUT-based impelementation example

\pgfmathsetmacro{\lutBestErrM}{1000}
\pgfmathsetmacro{\lutyM}{0.37}
\foreach \k in {0,...,1000}{%
        \pgfmathsetmacro{\lutyTmp}{\k/1000}%
        \pgfmathsetmacro{\lutErrTmp}{abs(xcurrentnorm(\lutyTmp)-xmtpvnorm(\lutyTmp))}%
        \pgfmathparse{\lutErrTmp < \lutBestErrM ? 1 : 0}%
        \ifnum\pgfmathresult=1
            \xdef\lutBestErrM{\lutErrTmp}%
            \xdef\lutyM{\lutyTmp}%
        \fi
    }
\pgfmathsetmacro{\lutxM}{xmtpvnorm(\lutyM)}

\pgfmathsetmacro{\lutBestErrR}{1000}
\pgfmathsetmacro{\lutyR}{0.81}
\foreach \k in {0,...,1000}{%
        \pgfmathsetmacro{\lutyTmp}{\k/1000}%
        \pgfmathsetmacro{\lutErrTmp}{abs(xcurrentnorm(\lutyTmp)-xmtpcnorm(\lutyTmp))}%
        \pgfmathparse{\lutErrTmp < \lutBestErrR ? 1 : 0}%
        \ifnum\pgfmathresult=1
            \xdef\lutBestErrR{\lutErrTmp}%
            \xdef\lutyR{\lutyTmp}%
        \fi
    }
\pgfmathsetmacro{\lutxR}{xmtpcnorm(\lutyR)}

\pgfmathsetmacro{\lutPsiN}{sqrt((\psim+\Ld*\Imax*\lutxR)^2+(\Lq*\Imax*\lutyR)^2)}
\pgfmathsetmacro{\lutTN}{torque(\Imax*\lutxR,\Imax*\lutyR)}

% A non-arbitrary teaching point for the T-psi-to-current map:
% midpoint between MTPV-current and MTPC-current intersections.
\pgfmathsetmacro{\lutyF}{0.5*(\lutyM+\lutyR)}
\pgfmathsetmacro{\lutxF}{xcurrentnorm(\lutyF)}
\pgfmathsetmacro{\lutPsiF}{sqrt((\psim+\Ld*\Imax*\lutxF)^2+(\Lq*\Imax*\lutyF)^2)}
\pgfmathsetmacro{\lutTF}{torque(\Imax*\lutxF,\Imax*\lutyF)}

\pgfmathdeclarefunction{lutfluxnorm}{2}{%
    \pgfmathparse{sqrt((\psim+\Ld*\Imax*(#1))^2+(\Lq*\Imax*(#2))^2)/\lutPsiN}%
}
\pgfmathdeclarefunction{luttorquenorm}{2}{%
    \pgfmathparse{torque(\Imax*(#1),\Imax*(#2))/\lutTN}%
}
\pgfmathdeclarefunction{lutfluxraw}{2}{%
    \pgfmathparse{sqrt((\psim+\Ld*\Imax*(#1))^2+(\Lq*\Imax*(#2))^2)}%
}